% Your own text at the start of the section "Two pawns vs one".
% make_book never writes this file.

\paragraph{Queen against pawn}\index{King and pawn endings!Queen against pawn}
Many races in pawn endings end with one side queening while the other side's pawn is one step
from its own queening square. The result then depends on the pawn. Against a central pawn or a
knight's pawn, the queen wins. A bishop's pawn or a rook's pawn may still have a chance, thanks
to stalemate. For the winning method and the exceptions, see Dvoretsky's Endgame Manual
\cite{Dvoretsky:2008-endgame-manual}.

\paragraph{Bähr's rule}\label{bahr-rule}\index{King and pawn endings!Bähr's rule}
Two rook's pawns often block each other while one side has a passed pawn on the other wing. The
stronger side's king goes after the enemy rook's pawn; the defender takes the passed pawn and runs
back to the corner in front of his own rook's pawn. Who comes first? Bähr gave a rule for this in
1936; Dvoretsky offers a handier version~\cite{Dvoretsky:2008-endgame-manual}. Here it is with
White as the stronger side, the rook's pawns on the a-file and the passed pawn on the king side.
\begin{itemize}
\item If White's rook's pawn has crossed the middle of the board, he always wins.
\item The \emph{normal} position is a draw. In it, the rook's pawns stand on either side of the
middle (a4 against a5), and the passed pawn has reached or crossed the key diagonal h3--c8, so the
black king reaches c8 in time.
\item Every square the passed pawn stands behind this diagonal is an extra tempo for White, and so
is his king standing in front of the passed pawn instead of beside it. Every square the rook's
pawns stand further back than in the normal position is an extra tempo for Black.
\item White wins only if he has more tempi than Black.
\end{itemize}
